% =====================================================================
% CHAPTER 3  DESIGN
% Smart-Edge Accessible Campus Navigation System
% Revised chapter. Replace the Design chapter of the report with this file
% (for example with % =====================================================================
% CHAPTER 3  DESIGN
% Smart-Edge Accessible Campus Navigation System
% Revised chapter. Replace the Design chapter of the report with this file
% (for example with % =====================================================================
% CHAPTER 3  DESIGN
% Smart-Edge Accessible Campus Navigation System
% Revised chapter. Replace the Design chapter of the report with this file
% (for example with % =====================================================================
% CHAPTER 3  DESIGN
% Smart-Edge Accessible Campus Navigation System
% Revised chapter. Replace the Design chapter of the report with this file
% (for example with \input{Chapter3_Design}). Figures are in ./figures/.
% Black text is the team's original text.
% Blue text is an update by Dr. Sekmen.
% Red text is a note from Dr. Sekmen and must be removed when addressed.
% =====================================================================
\providecommand{\upd}[1]{\textcolor{blue}{#1}}
\providecommand{\note}[1]{\textcolor{red}{#1}}
\providecommand{\fignote}[1]{{\color{red}\noindent\textbf{Note from Dr.\ Sekmen.} Figure~\ref{#1} shows suggested components for this subsystem. Check it, update it if needed, and explain each component.\par}}
\providecommand{\eqnote}{{\color{red}\noindent\textbf{Note from Dr.\ Sekmen.} This equation is a suggestion. Check it, and keep it with your own explanation only if you use this method.\par}}

\chapter{Design}
This chapter explains the design of our Smart-Edge Accessible Campus Navigation System for blind and low-vision pedestrians. The system uses an iPhone, Meta smart glasses, and campus map data to help users navigate around Tennessee State University. It is made up of several subsystems that work together to track the user's location, find accessible routes, recognize landmarks, and provide navigation instructions.

Section~\ref{sec:architecture} explains the overall system architecture and how the main components connect. Section~\ref{sec:detailed} explains each subsystem in more detail, including how information is collected, processed, and used during navigation. \upd{The figures in this chapter were drawn with assistance from Claude.}

% ---------------------------------------------------------------------
\section{System Architecture}
\label{sec:architecture}

Our Smart-Edge Accessible Campus Navigation System is designed to help blind and low-vision pedestrians move around campus using an iPhone and connected devices. The system is divided into several subsystems responsible for campus mapping, location tracking, route planning, landmark recognition, and navigation guidance. Each subsystem handles a different task, but they work together to guide the user to their destination.

The iPhone will serve as the main device responsible for processing information. It will use its built-in sensors, including GPS, compass, accelerometer, and gyroscope, to track the user's location and movement. The application will also use campus map data containing sidewalks, crossings, building entrances, and other important locations. This information will help the system determine where the user is and find a suitable route to their destination.

The Meta smart glasses will provide additional features through their camera, microphone, and speakers. Where supported, the camera will capture images of the user's surroundings, which will be processed on the iPhone to identify selected landmarks. The microphone will allow the user to give supported voice commands, while the speakers can provide spoken directions. If the glasses are disconnected, the application will still support navigation through the iPhone.

Before using the application, our team will collect information about selected campus routes and prepare the data needed for navigation. We will also train and test a computer vision model using a laptop. After preparing the route data and training the model, we will add them to the iPhone application. During navigation, the application will use this information along with live sensor readings to track the user's movement and provide directions.

The system will communicate with the user mainly through spoken instructions and vibration alerts. These instructions will help users identify upcoming turns, crossings, route changes, and when they have reached their destination. An Apple Watch may also be included as an optional device to provide additional vibration feedback and basic navigation controls.

Figure~\ref{fig:architecture} \upd{shows the} overall architecture of the system and how the different devices and subsystems work together. The design separates data preparation, navigation processing, user interaction, and testing so that each part can be developed and tested individually. Most of the navigation processing will take place on the iPhone, while the connected devices will provide additional ways for users to interact with the application.

\begin{figure}[htbp]
	\centering
	\includegraphics[width=\textwidth]{./figures/SA_P4_CampusNav.png}
	\caption[Overall system architecture of the Smart-Edge navigation system.]{Overall system architecture of the Smart-Edge navigation system.}
	\label{fig:architecture}
\end{figure}

% ---------------------------------------------------------------------
\section{Detailed Design}
\label{sec:detailed}

This section explains how the different subsystems of our navigation system will work. It builds on the architecture introduced in Section~\ref{sec:architecture} and provides more details about how the system will support accessible navigation around Tennessee State University. The design follows the functional requirements discussed in Chapter~2.

The iPhone will handle the main navigation tasks, including location tracking, route planning, and spoken guidance. Campus mapping and model training will be done before the application is deployed. Meta smart glasses will provide additional camera and audio features where supported, while an optional Apple Watch may be used for vibration alerts and basic controls.

{\color{red}\noindent\textbf{Note from Dr.\ Sekmen.} The detailed design should go one level below Figure~\ref{fig:architecture}. I gave each subsystem its own subsection with the same name as in the figure, moved your paragraphs under them, and added a figure with suggested components for each subsystem. In each figure the components are numbered in the order the data moves, inputs come in from the left and outputs leave on the right. For each subsystem, describe every component in a short paragraph that starts with its name in bold.\par}

% .....................................................................
\subsection[Campus Mapping]{\upd{Campus Mapping}}

\begin{figure}[htbp]
	\centering
	\includegraphics[width=0.9\textwidth]{./figures/Sub_CampusMapping.png}
	\caption[Campus Mapping and its components.]{\upd{Campus Mapping and its components.}}
	\label{fig:sub-map}
\end{figure}

\fignote{fig:sub-map}

The first step in developing our navigation system is collecting and preparing the campus data. Our team will map selected pedestrian routes around Tennessee State University and prepare the images needed to train the landmark recognition model. This work will be done before the navigation application is deployed.

For campus mapping, we will use OpenStreetMap, available TSU campus maps, and information collected by physically walking through the selected routes. We will identify important features such as sidewalks, intersections, crossings, ramps, stairs, building entrances, and landmarks. The routes will be organized into a pedestrian graph made up of nodes and edges. Nodes will represent important locations or points where the user may need to make a decision, while edges will represent the walkable paths connecting them. We will also record location and accessibility information for these features.

The collected route information will be stored using GeoJSON, a file format used to represent geographic data. This will allow us to store coordinates and information about each route in a structured format. We will also collect photographs of selected landmarks, such as building signs, entrances, and crosswalks. These images will be organized and used to train the landmark recognition model.

% .....................................................................
\subsection[Landmark Model Training]{\upd{Landmark Model Training}}

\begin{figure}[htbp]
	\centering
	\includegraphics[width=0.9\textwidth]{./figures/Sub_ModelTraining.png}
	\caption[Landmark Model Training and its components.]{\upd{Landmark Model Training and its components.}}
	\label{fig:sub-train}
\end{figure}

\fignote{fig:sub-train}

Our team will train or fine-tune a small computer vision model to recognize the selected landmarks. After training and testing the model, we plan to convert it into a format compatible with Apple's Core ML framework so it can run on the iPhone. The final route data and trained model will then be included in the application. This means the system will not need to train a model each time someone uses it.

% .....................................................................
\subsection[Positioning]{\upd{Positioning}}

\begin{figure}[htbp]
	\centering
	\includegraphics[width=0.9\textwidth]{./figures/Sub_Positioning.png}
	\caption[Positioning and its components.]{\upd{Positioning and its components.}}
	\label{fig:sub-pos}
\end{figure}

\fignote{fig:sub-pos}

The Positioning and Route Planning subsystems will work together to track the user's location and find an accessible route to their destination.

The Positioning subsystem will collect location and compass information using Apple's Core Location framework. It will also use Core Motion to collect data from the accelerometer and gyroscope. These sensors will help the system estimate the user's location, direction, and movement. Since GPS readings are not always accurate, especially near buildings, we plan to use map matching to compare the estimated location with the mapped pedestrian paths. This will help the application determine which walkway the user is most likely on.

\upd{For map matching, the closest point to the position $\mathbf{p}$ on an edge from node $\mathbf{a}$ to node $\mathbf{b}$, and its distance $d$, are
\begin{equation}
	t = \mathrm{clip}\!\left(\frac{(\mathbf{p}-\mathbf{a})\cdot(\mathbf{b}-\mathbf{a})}{\lVert \mathbf{b}-\mathbf{a} \rVert^2},\ 0,\ 1\right), \qquad
	d = \lVert \mathbf{p} - (\mathbf{a} + t(\mathbf{b}-\mathbf{a})) \rVert
\end{equation}
The user is placed on the edge with the smallest $d$.}

\eqnote

% .....................................................................
\subsection[Route Planning]{\upd{Route Planning}}

\begin{figure}[htbp]
	\centering
	\includegraphics[width=0.9\textwidth]{./figures/Sub_RoutePlanning.png}
	\caption[Route Planning and its components.]{\upd{Route Planning and its components.}}
	\label{fig:sub-route}
\end{figure}

\fignote{fig:sub-route}

The Route Planning subsystem will use the campus pedestrian graph stored in the application. When the user selects a destination, the system will use their estimated current location and a pathfinding algorithm, such as Dijkstra's algorithm or A*, to calculate a route. The selected algorithm will search the available paths to find a suitable way to reach the destination.

When choosing a route, the system will consider both distance and the accessibility information collected during campus mapping. For example, if a user requires a route without stairs, the system will avoid paths marked as containing stairs. Once a route has been calculated, the application will identify important points along the path, including turns and intersections, and pass this information to the navigation guidance subsystem.

\upd{Each edge $e$ can be given a cost
\begin{equation}
	c(e) =
	\begin{cases}
		\infty & \text{if } e \text{ has stairs and the user chose to avoid stairs} \\
		L(e)\, m(e) & \text{otherwise}
	\end{cases}
\end{equation}
where $L(e)$ is the length of the edge in meters and $m(e) \geq 1$ is a penalty that is 1 for a smooth accessible sidewalk and larger for a rough surface, a steep slope or a street crossing.}

\eqnote

As the user moves, the Positioning subsystem will continue updating their estimated location. If the system detects that the user has moved away from the planned route, it will attempt to calculate another accessible path. However, if the location information is unreliable or another suitable route cannot be found, the application will notify the user instead of giving directions that may not be correct.

{\color{red}\noindent\textbf{Note from Dr.\ Sekmen.} Graph loading and graph validation (R2) are not described yet. Also choose one search algorithm, and give the distance at which the user is considered off the route.\par}

% .....................................................................
\subsection[Accessible Interface and Guidance]{\upd{Accessible Interface and Guidance}}

\begin{figure}[htbp]
	\centering
	\includegraphics[width=0.9\textwidth]{./figures/Sub_Guidance.png}
	\caption[Accessible Interface and Guidance and its components.]{\upd{Accessible Interface and Guidance and its components.}}
	\label{fig:sub-guide}
\end{figure}

\fignote{fig:sub-guide}

The \upd{Accessible Interface and Guidance} subsystem will control how users interact with the application. Since the system is designed for blind and low-vision pedestrians, our main focus is to make navigation possible without requiring users to constantly look at their phones.

The iPhone application will allow users to select a destination, start or stop navigation, repeat instructions, and request their current location. We will also design the interface to work with Apple's VoiceOver accessibility feature, which reads information on the screen aloud. Supported voice commands will give users another way to perform common actions without using the touchscreen.

During navigation, this subsystem will receive information about the planned route, upcoming turns, and the user's estimated location. It will use this information to decide when to give directions. For example, when a user is approaching an intersection, the application will announce the upcoming turn before they reach it. This should give the user enough time to prepare for the next direction.

The system will use the iPhone's text-to-speech capabilities to provide spoken instructions. Vibration alerts will also be used for events such as turns, route deviations, warnings, and arrival at the destination. Users will be able to repeat instructions or request an update about their estimated location when needed.

We also plan to explore Apple Watch integration as an optional feature. If implemented, the watch could provide vibration alerts and simple controls, such as repeating an instruction or stopping navigation. The iPhone will still handle the main navigation processing, so users will not need an Apple Watch to use the system.

{\color{red}\noindent\textbf{Note from Dr.\ Sekmen.} Give the distance before a turn at which the instruction is spoken, and list the vibration pattern for each event in a small table.\par}

% .....................................................................
\subsection[Landmark Recognition]{\upd{Landmark Recognition}}

\begin{figure}[htbp]
	\centering
	\includegraphics[width=0.9\textwidth]{./figures/Sub_LandmarkRecognition.png}
	\caption[Landmark Recognition and its components.]{\upd{Landmark Recognition and its components.}}
	\label{fig:sub-recog}
\end{figure}

\fignote{fig:sub-recog}

For landmark recognition, camera images will be processed using the trained Core ML model stored in the application. The model will attempt to identify selected landmarks, such as building signs or entrances. When a landmark is detected, the system will check the model's confidence level and compare the result with the user's estimated location.

For example, if the camera identifies a building sign, the application will check whether that building is located near the user. If the identification is reliable and matches the location data, the application can announce the landmark to the user. This extra check is important because the model may sometimes recognize an object incorrectly.

% .....................................................................
\subsection[Glasses Connection]{\upd{Glasses Connection}}

\begin{figure}[htbp]
	\centering
	\includegraphics[width=0.9\textwidth]{./figures/Sub_GlassesConnection.png}
	\caption[Glasses Connection and its components.]{\upd{Glasses Connection and its components.}}
	\label{fig:sub-glasses}
\end{figure}

\fignote{fig:sub-glasses}

The Smart Glasses Integration and Landmark Recognition subsystems will provide additional ways for users to interact with the system and receive information about their surroundings. We plan to use compatible Meta smart glasses because they include a camera, microphone, and speakers. The iPhone will remain responsible for processing navigation information.

The Glasses Connection subsystem will manage communication between the smart glasses and the iPhone using supported Meta developer tools. Depending on the features available to developers, the glasses may allow the application to receive camera images, accept voice input, and play navigation instructions through their speakers. We will need to confirm which of these features are supported during implementation.

The Glasses Connection subsystem will also monitor whether the smart glasses are connected. If the connection is lost, the user will still be able to continue navigation using the iPhone. Where supported, the iPhone's camera, microphone, and speakers can provide some of the same functions. This will allow the main navigation features to work without depending entirely on the smart glasses.

{\color{red}\noindent\textbf{Note from Dr.\ Sekmen.} Confirm now which Meta toolkit you will use and which features it gives you (camera frames, microphone, speaker). If camera frames are not available, say that Landmark Recognition will use the phone camera.\par}

% .....................................................................
\subsection[Testing and Replay]{\upd{Testing and Replay}}

\begin{figure}[htbp]
	\centering
	\includegraphics[width=0.9\textwidth]{./figures/Sub_TestingReplay.png}
	\caption[Testing and Replay and its components.]{\upd{Testing and Replay and its components.}}
	\label{fig:sub-test}
\end{figure}

\fignote{fig:sub-test}

{\color{red}\noindent\textbf{Note from Dr.\ Sekmen.} This subsystem (R10 to R12) is in Figure~\ref{fig:architecture} but is not described in your chapter. Write this subsection. Explain how a test walk is recorded, how the recording is replayed through the app, how the replayed events are compared with the expected ones, how the supervised walking trials are run and what is measured, and how the accessibility review is done.\par}
). Figures are in ./figures/.
% Black text is the team's original text.
% Blue text is an update by Dr. Sekmen.
% Red text is a note from Dr. Sekmen and must be removed when addressed.
% =====================================================================
\providecommand{\upd}[1]{\textcolor{blue}{#1}}
\providecommand{\note}[1]{\textcolor{red}{#1}}
\providecommand{\fignote}[1]{{\color{red}\noindent\textbf{Note from Dr.\ Sekmen.} Figure~\ref{#1} shows suggested components for this subsystem. Check it, update it if needed, and explain each component.\par}}
\providecommand{\eqnote}{{\color{red}\noindent\textbf{Note from Dr.\ Sekmen.} This equation is a suggestion. Check it, and keep it with your own explanation only if you use this method.\par}}

\chapter{Design}
This chapter explains the design of our Smart-Edge Accessible Campus Navigation System for blind and low-vision pedestrians. The system uses an iPhone, Meta smart glasses, and campus map data to help users navigate around Tennessee State University. It is made up of several subsystems that work together to track the user's location, find accessible routes, recognize landmarks, and provide navigation instructions.

Section~\ref{sec:architecture} explains the overall system architecture and how the main components connect. Section~\ref{sec:detailed} explains each subsystem in more detail, including how information is collected, processed, and used during navigation. \upd{The figures in this chapter were drawn with assistance from Claude.}

% ---------------------------------------------------------------------
\section{System Architecture}
\label{sec:architecture}

Our Smart-Edge Accessible Campus Navigation System is designed to help blind and low-vision pedestrians move around campus using an iPhone and connected devices. The system is divided into several subsystems responsible for campus mapping, location tracking, route planning, landmark recognition, and navigation guidance. Each subsystem handles a different task, but they work together to guide the user to their destination.

The iPhone will serve as the main device responsible for processing information. It will use its built-in sensors, including GPS, compass, accelerometer, and gyroscope, to track the user's location and movement. The application will also use campus map data containing sidewalks, crossings, building entrances, and other important locations. This information will help the system determine where the user is and find a suitable route to their destination.

The Meta smart glasses will provide additional features through their camera, microphone, and speakers. Where supported, the camera will capture images of the user's surroundings, which will be processed on the iPhone to identify selected landmarks. The microphone will allow the user to give supported voice commands, while the speakers can provide spoken directions. If the glasses are disconnected, the application will still support navigation through the iPhone.

Before using the application, our team will collect information about selected campus routes and prepare the data needed for navigation. We will also train and test a computer vision model using a laptop. After preparing the route data and training the model, we will add them to the iPhone application. During navigation, the application will use this information along with live sensor readings to track the user's movement and provide directions.

The system will communicate with the user mainly through spoken instructions and vibration alerts. These instructions will help users identify upcoming turns, crossings, route changes, and when they have reached their destination. An Apple Watch may also be included as an optional device to provide additional vibration feedback and basic navigation controls.

Figure~\ref{fig:architecture} \upd{shows the} overall architecture of the system and how the different devices and subsystems work together. The design separates data preparation, navigation processing, user interaction, and testing so that each part can be developed and tested individually. Most of the navigation processing will take place on the iPhone, while the connected devices will provide additional ways for users to interact with the application.

\begin{figure}[htbp]
	\centering
	\includegraphics[width=\textwidth]{./figures/SA_P4_CampusNav.png}
	\caption[Overall system architecture of the Smart-Edge navigation system.]{Overall system architecture of the Smart-Edge navigation system.}
	\label{fig:architecture}
\end{figure}

% ---------------------------------------------------------------------
\section{Detailed Design}
\label{sec:detailed}

This section explains how the different subsystems of our navigation system will work. It builds on the architecture introduced in Section~\ref{sec:architecture} and provides more details about how the system will support accessible navigation around Tennessee State University. The design follows the functional requirements discussed in Chapter~2.

The iPhone will handle the main navigation tasks, including location tracking, route planning, and spoken guidance. Campus mapping and model training will be done before the application is deployed. Meta smart glasses will provide additional camera and audio features where supported, while an optional Apple Watch may be used for vibration alerts and basic controls.

{\color{red}\noindent\textbf{Note from Dr.\ Sekmen.} The detailed design should go one level below Figure~\ref{fig:architecture}. I gave each subsystem its own subsection with the same name as in the figure, moved your paragraphs under them, and added a figure with suggested components for each subsystem. In each figure the components are numbered in the order the data moves, inputs come in from the left and outputs leave on the right. For each subsystem, describe every component in a short paragraph that starts with its name in bold.\par}

% .....................................................................
\subsection[Campus Mapping]{\upd{Campus Mapping}}

\begin{figure}[htbp]
	\centering
	\includegraphics[width=0.9\textwidth]{./figures/Sub_CampusMapping.png}
	\caption[Campus Mapping and its components.]{\upd{Campus Mapping and its components.}}
	\label{fig:sub-map}
\end{figure}

\fignote{fig:sub-map}

The first step in developing our navigation system is collecting and preparing the campus data. Our team will map selected pedestrian routes around Tennessee State University and prepare the images needed to train the landmark recognition model. This work will be done before the navigation application is deployed.

For campus mapping, we will use OpenStreetMap, available TSU campus maps, and information collected by physically walking through the selected routes. We will identify important features such as sidewalks, intersections, crossings, ramps, stairs, building entrances, and landmarks. The routes will be organized into a pedestrian graph made up of nodes and edges. Nodes will represent important locations or points where the user may need to make a decision, while edges will represent the walkable paths connecting them. We will also record location and accessibility information for these features.

The collected route information will be stored using GeoJSON, a file format used to represent geographic data. This will allow us to store coordinates and information about each route in a structured format. We will also collect photographs of selected landmarks, such as building signs, entrances, and crosswalks. These images will be organized and used to train the landmark recognition model.

% .....................................................................
\subsection[Landmark Model Training]{\upd{Landmark Model Training}}

\begin{figure}[htbp]
	\centering
	\includegraphics[width=0.9\textwidth]{./figures/Sub_ModelTraining.png}
	\caption[Landmark Model Training and its components.]{\upd{Landmark Model Training and its components.}}
	\label{fig:sub-train}
\end{figure}

\fignote{fig:sub-train}

Our team will train or fine-tune a small computer vision model to recognize the selected landmarks. After training and testing the model, we plan to convert it into a format compatible with Apple's Core ML framework so it can run on the iPhone. The final route data and trained model will then be included in the application. This means the system will not need to train a model each time someone uses it.

% .....................................................................
\subsection[Positioning]{\upd{Positioning}}

\begin{figure}[htbp]
	\centering
	\includegraphics[width=0.9\textwidth]{./figures/Sub_Positioning.png}
	\caption[Positioning and its components.]{\upd{Positioning and its components.}}
	\label{fig:sub-pos}
\end{figure}

\fignote{fig:sub-pos}

The Positioning and Route Planning subsystems will work together to track the user's location and find an accessible route to their destination.

The Positioning subsystem will collect location and compass information using Apple's Core Location framework. It will also use Core Motion to collect data from the accelerometer and gyroscope. These sensors will help the system estimate the user's location, direction, and movement. Since GPS readings are not always accurate, especially near buildings, we plan to use map matching to compare the estimated location with the mapped pedestrian paths. This will help the application determine which walkway the user is most likely on.

\upd{For map matching, the closest point to the position $\mathbf{p}$ on an edge from node $\mathbf{a}$ to node $\mathbf{b}$, and its distance $d$, are
\begin{equation}
	t = \mathrm{clip}\!\left(\frac{(\mathbf{p}-\mathbf{a})\cdot(\mathbf{b}-\mathbf{a})}{\lVert \mathbf{b}-\mathbf{a} \rVert^2},\ 0,\ 1\right), \qquad
	d = \lVert \mathbf{p} - (\mathbf{a} + t(\mathbf{b}-\mathbf{a})) \rVert
\end{equation}
The user is placed on the edge with the smallest $d$.}

\eqnote

% .....................................................................
\subsection[Route Planning]{\upd{Route Planning}}

\begin{figure}[htbp]
	\centering
	\includegraphics[width=0.9\textwidth]{./figures/Sub_RoutePlanning.png}
	\caption[Route Planning and its components.]{\upd{Route Planning and its components.}}
	\label{fig:sub-route}
\end{figure}

\fignote{fig:sub-route}

The Route Planning subsystem will use the campus pedestrian graph stored in the application. When the user selects a destination, the system will use their estimated current location and a pathfinding algorithm, such as Dijkstra's algorithm or A*, to calculate a route. The selected algorithm will search the available paths to find a suitable way to reach the destination.

When choosing a route, the system will consider both distance and the accessibility information collected during campus mapping. For example, if a user requires a route without stairs, the system will avoid paths marked as containing stairs. Once a route has been calculated, the application will identify important points along the path, including turns and intersections, and pass this information to the navigation guidance subsystem.

\upd{Each edge $e$ can be given a cost
\begin{equation}
	c(e) =
	\begin{cases}
		\infty & \text{if } e \text{ has stairs and the user chose to avoid stairs} \\
		L(e)\, m(e) & \text{otherwise}
	\end{cases}
\end{equation}
where $L(e)$ is the length of the edge in meters and $m(e) \geq 1$ is a penalty that is 1 for a smooth accessible sidewalk and larger for a rough surface, a steep slope or a street crossing.}

\eqnote

As the user moves, the Positioning subsystem will continue updating their estimated location. If the system detects that the user has moved away from the planned route, it will attempt to calculate another accessible path. However, if the location information is unreliable or another suitable route cannot be found, the application will notify the user instead of giving directions that may not be correct.

{\color{red}\noindent\textbf{Note from Dr.\ Sekmen.} Graph loading and graph validation (R2) are not described yet. Also choose one search algorithm, and give the distance at which the user is considered off the route.\par}

% .....................................................................
\subsection[Accessible Interface and Guidance]{\upd{Accessible Interface and Guidance}}

\begin{figure}[htbp]
	\centering
	\includegraphics[width=0.9\textwidth]{./figures/Sub_Guidance.png}
	\caption[Accessible Interface and Guidance and its components.]{\upd{Accessible Interface and Guidance and its components.}}
	\label{fig:sub-guide}
\end{figure}

\fignote{fig:sub-guide}

The \upd{Accessible Interface and Guidance} subsystem will control how users interact with the application. Since the system is designed for blind and low-vision pedestrians, our main focus is to make navigation possible without requiring users to constantly look at their phones.

The iPhone application will allow users to select a destination, start or stop navigation, repeat instructions, and request their current location. We will also design the interface to work with Apple's VoiceOver accessibility feature, which reads information on the screen aloud. Supported voice commands will give users another way to perform common actions without using the touchscreen.

During navigation, this subsystem will receive information about the planned route, upcoming turns, and the user's estimated location. It will use this information to decide when to give directions. For example, when a user is approaching an intersection, the application will announce the upcoming turn before they reach it. This should give the user enough time to prepare for the next direction.

The system will use the iPhone's text-to-speech capabilities to provide spoken instructions. Vibration alerts will also be used for events such as turns, route deviations, warnings, and arrival at the destination. Users will be able to repeat instructions or request an update about their estimated location when needed.

We also plan to explore Apple Watch integration as an optional feature. If implemented, the watch could provide vibration alerts and simple controls, such as repeating an instruction or stopping navigation. The iPhone will still handle the main navigation processing, so users will not need an Apple Watch to use the system.

{\color{red}\noindent\textbf{Note from Dr.\ Sekmen.} Give the distance before a turn at which the instruction is spoken, and list the vibration pattern for each event in a small table.\par}

% .....................................................................
\subsection[Landmark Recognition]{\upd{Landmark Recognition}}

\begin{figure}[htbp]
	\centering
	\includegraphics[width=0.9\textwidth]{./figures/Sub_LandmarkRecognition.png}
	\caption[Landmark Recognition and its components.]{\upd{Landmark Recognition and its components.}}
	\label{fig:sub-recog}
\end{figure}

\fignote{fig:sub-recog}

For landmark recognition, camera images will be processed using the trained Core ML model stored in the application. The model will attempt to identify selected landmarks, such as building signs or entrances. When a landmark is detected, the system will check the model's confidence level and compare the result with the user's estimated location.

For example, if the camera identifies a building sign, the application will check whether that building is located near the user. If the identification is reliable and matches the location data, the application can announce the landmark to the user. This extra check is important because the model may sometimes recognize an object incorrectly.

% .....................................................................
\subsection[Glasses Connection]{\upd{Glasses Connection}}

\begin{figure}[htbp]
	\centering
	\includegraphics[width=0.9\textwidth]{./figures/Sub_GlassesConnection.png}
	\caption[Glasses Connection and its components.]{\upd{Glasses Connection and its components.}}
	\label{fig:sub-glasses}
\end{figure}

\fignote{fig:sub-glasses}

The Smart Glasses Integration and Landmark Recognition subsystems will provide additional ways for users to interact with the system and receive information about their surroundings. We plan to use compatible Meta smart glasses because they include a camera, microphone, and speakers. The iPhone will remain responsible for processing navigation information.

The Glasses Connection subsystem will manage communication between the smart glasses and the iPhone using supported Meta developer tools. Depending on the features available to developers, the glasses may allow the application to receive camera images, accept voice input, and play navigation instructions through their speakers. We will need to confirm which of these features are supported during implementation.

The Glasses Connection subsystem will also monitor whether the smart glasses are connected. If the connection is lost, the user will still be able to continue navigation using the iPhone. Where supported, the iPhone's camera, microphone, and speakers can provide some of the same functions. This will allow the main navigation features to work without depending entirely on the smart glasses.

{\color{red}\noindent\textbf{Note from Dr.\ Sekmen.} Confirm now which Meta toolkit you will use and which features it gives you (camera frames, microphone, speaker). If camera frames are not available, say that Landmark Recognition will use the phone camera.\par}

% .....................................................................
\subsection[Testing and Replay]{\upd{Testing and Replay}}

\begin{figure}[htbp]
	\centering
	\includegraphics[width=0.9\textwidth]{./figures/Sub_TestingReplay.png}
	\caption[Testing and Replay and its components.]{\upd{Testing and Replay and its components.}}
	\label{fig:sub-test}
\end{figure}

\fignote{fig:sub-test}

{\color{red}\noindent\textbf{Note from Dr.\ Sekmen.} This subsystem (R10 to R12) is in Figure~\ref{fig:architecture} but is not described in your chapter. Write this subsection. Explain how a test walk is recorded, how the recording is replayed through the app, how the replayed events are compared with the expected ones, how the supervised walking trials are run and what is measured, and how the accessibility review is done.\par}
). Figures are in ./figures/.
% Black text is the team's original text.
% Blue text is an update by Dr. Sekmen.
% Red text is a note from Dr. Sekmen and must be removed when addressed.
% =====================================================================
\providecommand{\upd}[1]{\textcolor{blue}{#1}}
\providecommand{\note}[1]{\textcolor{red}{#1}}
\providecommand{\fignote}[1]{{\color{red}\noindent\textbf{Note from Dr.\ Sekmen.} Figure~\ref{#1} shows suggested components for this subsystem. Check it, update it if needed, and explain each component.\par}}
\providecommand{\eqnote}{{\color{red}\noindent\textbf{Note from Dr.\ Sekmen.} This equation is a suggestion. Check it, and keep it with your own explanation only if you use this method.\par}}

\chapter{Design}
This chapter explains the design of our Smart-Edge Accessible Campus Navigation System for blind and low-vision pedestrians. The system uses an iPhone, Meta smart glasses, and campus map data to help users navigate around Tennessee State University. It is made up of several subsystems that work together to track the user's location, find accessible routes, recognize landmarks, and provide navigation instructions.

Section~\ref{sec:architecture} explains the overall system architecture and how the main components connect. Section~\ref{sec:detailed} explains each subsystem in more detail, including how information is collected, processed, and used during navigation. \upd{The figures in this chapter were drawn with assistance from Claude.}

% ---------------------------------------------------------------------
\section{System Architecture}
\label{sec:architecture}

Our Smart-Edge Accessible Campus Navigation System is designed to help blind and low-vision pedestrians move around campus using an iPhone and connected devices. The system is divided into several subsystems responsible for campus mapping, location tracking, route planning, landmark recognition, and navigation guidance. Each subsystem handles a different task, but they work together to guide the user to their destination.

The iPhone will serve as the main device responsible for processing information. It will use its built-in sensors, including GPS, compass, accelerometer, and gyroscope, to track the user's location and movement. The application will also use campus map data containing sidewalks, crossings, building entrances, and other important locations. This information will help the system determine where the user is and find a suitable route to their destination.

The Meta smart glasses will provide additional features through their camera, microphone, and speakers. Where supported, the camera will capture images of the user's surroundings, which will be processed on the iPhone to identify selected landmarks. The microphone will allow the user to give supported voice commands, while the speakers can provide spoken directions. If the glasses are disconnected, the application will still support navigation through the iPhone.

Before using the application, our team will collect information about selected campus routes and prepare the data needed for navigation. We will also train and test a computer vision model using a laptop. After preparing the route data and training the model, we will add them to the iPhone application. During navigation, the application will use this information along with live sensor readings to track the user's movement and provide directions.

The system will communicate with the user mainly through spoken instructions and vibration alerts. These instructions will help users identify upcoming turns, crossings, route changes, and when they have reached their destination. An Apple Watch may also be included as an optional device to provide additional vibration feedback and basic navigation controls.

Figure~\ref{fig:architecture} \upd{shows the} overall architecture of the system and how the different devices and subsystems work together. The design separates data preparation, navigation processing, user interaction, and testing so that each part can be developed and tested individually. Most of the navigation processing will take place on the iPhone, while the connected devices will provide additional ways for users to interact with the application.

\begin{figure}[htbp]
	\centering
	\includegraphics[width=\textwidth]{./figures/SA_P4_CampusNav.png}
	\caption[Overall system architecture of the Smart-Edge navigation system.]{Overall system architecture of the Smart-Edge navigation system.}
	\label{fig:architecture}
\end{figure}

% ---------------------------------------------------------------------
\section{Detailed Design}
\label{sec:detailed}

This section explains how the different subsystems of our navigation system will work. It builds on the architecture introduced in Section~\ref{sec:architecture} and provides more details about how the system will support accessible navigation around Tennessee State University. The design follows the functional requirements discussed in Chapter~2.

The iPhone will handle the main navigation tasks, including location tracking, route planning, and spoken guidance. Campus mapping and model training will be done before the application is deployed. Meta smart glasses will provide additional camera and audio features where supported, while an optional Apple Watch may be used for vibration alerts and basic controls.

{\color{red}\noindent\textbf{Note from Dr.\ Sekmen.} The detailed design should go one level below Figure~\ref{fig:architecture}. I gave each subsystem its own subsection with the same name as in the figure, moved your paragraphs under them, and added a figure with suggested components for each subsystem. In each figure the components are numbered in the order the data moves, inputs come in from the left and outputs leave on the right. For each subsystem, describe every component in a short paragraph that starts with its name in bold.\par}

% .....................................................................
\subsection[Campus Mapping]{\upd{Campus Mapping}}

\begin{figure}[htbp]
	\centering
	\includegraphics[width=0.9\textwidth]{./figures/Sub_CampusMapping.png}
	\caption[Campus Mapping and its components.]{\upd{Campus Mapping and its components.}}
	\label{fig:sub-map}
\end{figure}

\fignote{fig:sub-map}

The first step in developing our navigation system is collecting and preparing the campus data. Our team will map selected pedestrian routes around Tennessee State University and prepare the images needed to train the landmark recognition model. This work will be done before the navigation application is deployed.

For campus mapping, we will use OpenStreetMap, available TSU campus maps, and information collected by physically walking through the selected routes. We will identify important features such as sidewalks, intersections, crossings, ramps, stairs, building entrances, and landmarks. The routes will be organized into a pedestrian graph made up of nodes and edges. Nodes will represent important locations or points where the user may need to make a decision, while edges will represent the walkable paths connecting them. We will also record location and accessibility information for these features.

The collected route information will be stored using GeoJSON, a file format used to represent geographic data. This will allow us to store coordinates and information about each route in a structured format. We will also collect photographs of selected landmarks, such as building signs, entrances, and crosswalks. These images will be organized and used to train the landmark recognition model.

% .....................................................................
\subsection[Landmark Model Training]{\upd{Landmark Model Training}}

\begin{figure}[htbp]
	\centering
	\includegraphics[width=0.9\textwidth]{./figures/Sub_ModelTraining.png}
	\caption[Landmark Model Training and its components.]{\upd{Landmark Model Training and its components.}}
	\label{fig:sub-train}
\end{figure}

\fignote{fig:sub-train}

Our team will train or fine-tune a small computer vision model to recognize the selected landmarks. After training and testing the model, we plan to convert it into a format compatible with Apple's Core ML framework so it can run on the iPhone. The final route data and trained model will then be included in the application. This means the system will not need to train a model each time someone uses it.

% .....................................................................
\subsection[Positioning]{\upd{Positioning}}

\begin{figure}[htbp]
	\centering
	\includegraphics[width=0.9\textwidth]{./figures/Sub_Positioning.png}
	\caption[Positioning and its components.]{\upd{Positioning and its components.}}
	\label{fig:sub-pos}
\end{figure}

\fignote{fig:sub-pos}

The Positioning and Route Planning subsystems will work together to track the user's location and find an accessible route to their destination.

The Positioning subsystem will collect location and compass information using Apple's Core Location framework. It will also use Core Motion to collect data from the accelerometer and gyroscope. These sensors will help the system estimate the user's location, direction, and movement. Since GPS readings are not always accurate, especially near buildings, we plan to use map matching to compare the estimated location with the mapped pedestrian paths. This will help the application determine which walkway the user is most likely on.

\upd{For map matching, the closest point to the position $\mathbf{p}$ on an edge from node $\mathbf{a}$ to node $\mathbf{b}$, and its distance $d$, are
\begin{equation}
	t = \mathrm{clip}\!\left(\frac{(\mathbf{p}-\mathbf{a})\cdot(\mathbf{b}-\mathbf{a})}{\lVert \mathbf{b}-\mathbf{a} \rVert^2},\ 0,\ 1\right), \qquad
	d = \lVert \mathbf{p} - (\mathbf{a} + t(\mathbf{b}-\mathbf{a})) \rVert
\end{equation}
The user is placed on the edge with the smallest $d$.}

\eqnote

% .....................................................................
\subsection[Route Planning]{\upd{Route Planning}}

\begin{figure}[htbp]
	\centering
	\includegraphics[width=0.9\textwidth]{./figures/Sub_RoutePlanning.png}
	\caption[Route Planning and its components.]{\upd{Route Planning and its components.}}
	\label{fig:sub-route}
\end{figure}

\fignote{fig:sub-route}

The Route Planning subsystem will use the campus pedestrian graph stored in the application. When the user selects a destination, the system will use their estimated current location and a pathfinding algorithm, such as Dijkstra's algorithm or A*, to calculate a route. The selected algorithm will search the available paths to find a suitable way to reach the destination.

When choosing a route, the system will consider both distance and the accessibility information collected during campus mapping. For example, if a user requires a route without stairs, the system will avoid paths marked as containing stairs. Once a route has been calculated, the application will identify important points along the path, including turns and intersections, and pass this information to the navigation guidance subsystem.

\upd{Each edge $e$ can be given a cost
\begin{equation}
	c(e) =
	\begin{cases}
		\infty & \text{if } e \text{ has stairs and the user chose to avoid stairs} \\
		L(e)\, m(e) & \text{otherwise}
	\end{cases}
\end{equation}
where $L(e)$ is the length of the edge in meters and $m(e) \geq 1$ is a penalty that is 1 for a smooth accessible sidewalk and larger for a rough surface, a steep slope or a street crossing.}

\eqnote

As the user moves, the Positioning subsystem will continue updating their estimated location. If the system detects that the user has moved away from the planned route, it will attempt to calculate another accessible path. However, if the location information is unreliable or another suitable route cannot be found, the application will notify the user instead of giving directions that may not be correct.

{\color{red}\noindent\textbf{Note from Dr.\ Sekmen.} Graph loading and graph validation (R2) are not described yet. Also choose one search algorithm, and give the distance at which the user is considered off the route.\par}

% .....................................................................
\subsection[Accessible Interface and Guidance]{\upd{Accessible Interface and Guidance}}

\begin{figure}[htbp]
	\centering
	\includegraphics[width=0.9\textwidth]{./figures/Sub_Guidance.png}
	\caption[Accessible Interface and Guidance and its components.]{\upd{Accessible Interface and Guidance and its components.}}
	\label{fig:sub-guide}
\end{figure}

\fignote{fig:sub-guide}

The \upd{Accessible Interface and Guidance} subsystem will control how users interact with the application. Since the system is designed for blind and low-vision pedestrians, our main focus is to make navigation possible without requiring users to constantly look at their phones.

The iPhone application will allow users to select a destination, start or stop navigation, repeat instructions, and request their current location. We will also design the interface to work with Apple's VoiceOver accessibility feature, which reads information on the screen aloud. Supported voice commands will give users another way to perform common actions without using the touchscreen.

During navigation, this subsystem will receive information about the planned route, upcoming turns, and the user's estimated location. It will use this information to decide when to give directions. For example, when a user is approaching an intersection, the application will announce the upcoming turn before they reach it. This should give the user enough time to prepare for the next direction.

The system will use the iPhone's text-to-speech capabilities to provide spoken instructions. Vibration alerts will also be used for events such as turns, route deviations, warnings, and arrival at the destination. Users will be able to repeat instructions or request an update about their estimated location when needed.

We also plan to explore Apple Watch integration as an optional feature. If implemented, the watch could provide vibration alerts and simple controls, such as repeating an instruction or stopping navigation. The iPhone will still handle the main navigation processing, so users will not need an Apple Watch to use the system.

{\color{red}\noindent\textbf{Note from Dr.\ Sekmen.} Give the distance before a turn at which the instruction is spoken, and list the vibration pattern for each event in a small table.\par}

% .....................................................................
\subsection[Landmark Recognition]{\upd{Landmark Recognition}}

\begin{figure}[htbp]
	\centering
	\includegraphics[width=0.9\textwidth]{./figures/Sub_LandmarkRecognition.png}
	\caption[Landmark Recognition and its components.]{\upd{Landmark Recognition and its components.}}
	\label{fig:sub-recog}
\end{figure}

\fignote{fig:sub-recog}

For landmark recognition, camera images will be processed using the trained Core ML model stored in the application. The model will attempt to identify selected landmarks, such as building signs or entrances. When a landmark is detected, the system will check the model's confidence level and compare the result with the user's estimated location.

For example, if the camera identifies a building sign, the application will check whether that building is located near the user. If the identification is reliable and matches the location data, the application can announce the landmark to the user. This extra check is important because the model may sometimes recognize an object incorrectly.

% .....................................................................
\subsection[Glasses Connection]{\upd{Glasses Connection}}

\begin{figure}[htbp]
	\centering
	\includegraphics[width=0.9\textwidth]{./figures/Sub_GlassesConnection.png}
	\caption[Glasses Connection and its components.]{\upd{Glasses Connection and its components.}}
	\label{fig:sub-glasses}
\end{figure}

\fignote{fig:sub-glasses}

The Smart Glasses Integration and Landmark Recognition subsystems will provide additional ways for users to interact with the system and receive information about their surroundings. We plan to use compatible Meta smart glasses because they include a camera, microphone, and speakers. The iPhone will remain responsible for processing navigation information.

The Glasses Connection subsystem will manage communication between the smart glasses and the iPhone using supported Meta developer tools. Depending on the features available to developers, the glasses may allow the application to receive camera images, accept voice input, and play navigation instructions through their speakers. We will need to confirm which of these features are supported during implementation.

The Glasses Connection subsystem will also monitor whether the smart glasses are connected. If the connection is lost, the user will still be able to continue navigation using the iPhone. Where supported, the iPhone's camera, microphone, and speakers can provide some of the same functions. This will allow the main navigation features to work without depending entirely on the smart glasses.

{\color{red}\noindent\textbf{Note from Dr.\ Sekmen.} Confirm now which Meta toolkit you will use and which features it gives you (camera frames, microphone, speaker). If camera frames are not available, say that Landmark Recognition will use the phone camera.\par}

% .....................................................................
\subsection[Testing and Replay]{\upd{Testing and Replay}}

\begin{figure}[htbp]
	\centering
	\includegraphics[width=0.9\textwidth]{./figures/Sub_TestingReplay.png}
	\caption[Testing and Replay and its components.]{\upd{Testing and Replay and its components.}}
	\label{fig:sub-test}
\end{figure}

\fignote{fig:sub-test}

{\color{red}\noindent\textbf{Note from Dr.\ Sekmen.} This subsystem (R10 to R12) is in Figure~\ref{fig:architecture} but is not described in your chapter. Write this subsection. Explain how a test walk is recorded, how the recording is replayed through the app, how the replayed events are compared with the expected ones, how the supervised walking trials are run and what is measured, and how the accessibility review is done.\par}
). Figures are in ./figures/.
% Black text is the team's original text.
% Blue text is an update by Dr. Sekmen.
% Red text is a note from Dr. Sekmen and must be removed when addressed.
% =====================================================================
\providecommand{\upd}[1]{\textcolor{blue}{#1}}
\providecommand{\note}[1]{\textcolor{red}{#1}}
\providecommand{\fignote}[1]{{\color{red}\noindent\textbf{Note from Dr.\ Sekmen.} Figure~\ref{#1} shows suggested components for this subsystem. Check it, update it if needed, and explain each component.\par}}
\providecommand{\eqnote}{{\color{red}\noindent\textbf{Note from Dr.\ Sekmen.} This equation is a suggestion. Check it, and keep it with your own explanation only if you use this method.\par}}

\chapter{Design}
This chapter explains the design of our Smart-Edge Accessible Campus Navigation System for blind and low-vision pedestrians. The system uses an iPhone, Meta smart glasses, and campus map data to help users navigate around Tennessee State University. It is made up of several subsystems that work together to track the user's location, find accessible routes, recognize landmarks, and provide navigation instructions.

Section~\ref{sec:architecture} explains the overall system architecture and how the main components connect. Section~\ref{sec:detailed} explains each subsystem in more detail, including how information is collected, processed, and used during navigation. \upd{The figures in this chapter were drawn with assistance from Claude.}

% ---------------------------------------------------------------------
\section{System Architecture}
\label{sec:architecture}

Our Smart-Edge Accessible Campus Navigation System is designed to help blind and low-vision pedestrians move around campus using an iPhone and connected devices. The system is divided into several subsystems responsible for campus mapping, location tracking, route planning, landmark recognition, and navigation guidance. Each subsystem handles a different task, but they work together to guide the user to their destination.

The iPhone will serve as the main device responsible for processing information. It will use its built-in sensors, including GPS, compass, accelerometer, and gyroscope, to track the user's location and movement. The application will also use campus map data containing sidewalks, crossings, building entrances, and other important locations. This information will help the system determine where the user is and find a suitable route to their destination.

The Meta smart glasses will provide additional features through their camera, microphone, and speakers. Where supported, the camera will capture images of the user's surroundings, which will be processed on the iPhone to identify selected landmarks. The microphone will allow the user to give supported voice commands, while the speakers can provide spoken directions. If the glasses are disconnected, the application will still support navigation through the iPhone.

Before using the application, our team will collect information about selected campus routes and prepare the data needed for navigation. We will also train and test a computer vision model using a laptop. After preparing the route data and training the model, we will add them to the iPhone application. During navigation, the application will use this information along with live sensor readings to track the user's movement and provide directions.

The system will communicate with the user mainly through spoken instructions and vibration alerts. These instructions will help users identify upcoming turns, crossings, route changes, and when they have reached their destination. An Apple Watch may also be included as an optional device to provide additional vibration feedback and basic navigation controls.

Figure~\ref{fig:architecture} \upd{shows the} overall architecture of the system and how the different devices and subsystems work together. The design separates data preparation, navigation processing, user interaction, and testing so that each part can be developed and tested individually. Most of the navigation processing will take place on the iPhone, while the connected devices will provide additional ways for users to interact with the application.

\begin{figure}[htbp]
	\centering
	\includegraphics[width=\textwidth]{./figures/SA_P4_CampusNav.png}
	\caption[Overall system architecture of the Smart-Edge navigation system.]{Overall system architecture of the Smart-Edge navigation system.}
	\label{fig:architecture}
\end{figure}

% ---------------------------------------------------------------------
\section{Detailed Design}
\label{sec:detailed}

This section explains how the different subsystems of our navigation system will work. It builds on the architecture introduced in Section~\ref{sec:architecture} and provides more details about how the system will support accessible navigation around Tennessee State University. The design follows the functional requirements discussed in Chapter~2.

The iPhone will handle the main navigation tasks, including location tracking, route planning, and spoken guidance. Campus mapping and model training will be done before the application is deployed. Meta smart glasses will provide additional camera and audio features where supported, while an optional Apple Watch may be used for vibration alerts and basic controls.

{\color{red}\noindent\textbf{Note from Dr.\ Sekmen.} The detailed design should go one level below Figure~\ref{fig:architecture}. I gave each subsystem its own subsection with the same name as in the figure, moved your paragraphs under them, and added a figure with suggested components for each subsystem. In each figure the components are numbered in the order the data moves, inputs come in from the left and outputs leave on the right. For each subsystem, describe every component in a short paragraph that starts with its name in bold.\par}

% .....................................................................
\subsection[Campus Mapping]{\upd{Campus Mapping}}

\begin{figure}[htbp]
	\centering
	\includegraphics[width=0.9\textwidth]{./figures/Sub_CampusMapping.png}
	\caption[Campus Mapping and its components.]{\upd{Campus Mapping and its components.}}
	\label{fig:sub-map}
\end{figure}

\fignote{fig:sub-map}

The first step in developing our navigation system is collecting and preparing the campus data. Our team will map selected pedestrian routes around Tennessee State University and prepare the images needed to train the landmark recognition model. This work will be done before the navigation application is deployed.

For campus mapping, we will use OpenStreetMap, available TSU campus maps, and information collected by physically walking through the selected routes. We will identify important features such as sidewalks, intersections, crossings, ramps, stairs, building entrances, and landmarks. The routes will be organized into a pedestrian graph made up of nodes and edges. Nodes will represent important locations or points where the user may need to make a decision, while edges will represent the walkable paths connecting them. We will also record location and accessibility information for these features.

The collected route information will be stored using GeoJSON, a file format used to represent geographic data. This will allow us to store coordinates and information about each route in a structured format. We will also collect photographs of selected landmarks, such as building signs, entrances, and crosswalks. These images will be organized and used to train the landmark recognition model.

% .....................................................................
\subsection[Landmark Model Training]{\upd{Landmark Model Training}}

\begin{figure}[htbp]
	\centering
	\includegraphics[width=0.9\textwidth]{./figures/Sub_ModelTraining.png}
	\caption[Landmark Model Training and its components.]{\upd{Landmark Model Training and its components.}}
	\label{fig:sub-train}
\end{figure}

\fignote{fig:sub-train}

Our team will train or fine-tune a small computer vision model to recognize the selected landmarks. After training and testing the model, we plan to convert it into a format compatible with Apple's Core ML framework so it can run on the iPhone. The final route data and trained model will then be included in the application. This means the system will not need to train a model each time someone uses it.

% .....................................................................
\subsection[Positioning]{\upd{Positioning}}

\begin{figure}[htbp]
	\centering
	\includegraphics[width=0.9\textwidth]{./figures/Sub_Positioning.png}
	\caption[Positioning and its components.]{\upd{Positioning and its components.}}
	\label{fig:sub-pos}
\end{figure}

\fignote{fig:sub-pos}

The Positioning and Route Planning subsystems will work together to track the user's location and find an accessible route to their destination.

The Positioning subsystem will collect location and compass information using Apple's Core Location framework. It will also use Core Motion to collect data from the accelerometer and gyroscope. These sensors will help the system estimate the user's location, direction, and movement. Since GPS readings are not always accurate, especially near buildings, we plan to use map matching to compare the estimated location with the mapped pedestrian paths. This will help the application determine which walkway the user is most likely on.

\upd{For map matching, the closest point to the position $\mathbf{p}$ on an edge from node $\mathbf{a}$ to node $\mathbf{b}$, and its distance $d$, are
\begin{equation}
	t = \mathrm{clip}\!\left(\frac{(\mathbf{p}-\mathbf{a})\cdot(\mathbf{b}-\mathbf{a})}{\lVert \mathbf{b}-\mathbf{a} \rVert^2},\ 0,\ 1\right), \qquad
	d = \lVert \mathbf{p} - (\mathbf{a} + t(\mathbf{b}-\mathbf{a})) \rVert
\end{equation}
The user is placed on the edge with the smallest $d$.}

\eqnote

% .....................................................................
\subsection[Route Planning]{\upd{Route Planning}}

\begin{figure}[htbp]
	\centering
	\includegraphics[width=0.9\textwidth]{./figures/Sub_RoutePlanning.png}
	\caption[Route Planning and its components.]{\upd{Route Planning and its components.}}
	\label{fig:sub-route}
\end{figure}

\fignote{fig:sub-route}

The Route Planning subsystem will use the campus pedestrian graph stored in the application. When the user selects a destination, the system will use their estimated current location and a pathfinding algorithm, such as Dijkstra's algorithm or A*, to calculate a route. The selected algorithm will search the available paths to find a suitable way to reach the destination.

When choosing a route, the system will consider both distance and the accessibility information collected during campus mapping. For example, if a user requires a route without stairs, the system will avoid paths marked as containing stairs. Once a route has been calculated, the application will identify important points along the path, including turns and intersections, and pass this information to the navigation guidance subsystem.

\upd{Each edge $e$ can be given a cost
\begin{equation}
	c(e) =
	\begin{cases}
		\infty & \text{if } e \text{ has stairs and the user chose to avoid stairs} \\
		L(e)\, m(e) & \text{otherwise}
	\end{cases}
\end{equation}
where $L(e)$ is the length of the edge in meters and $m(e) \geq 1$ is a penalty that is 1 for a smooth accessible sidewalk and larger for a rough surface, a steep slope or a street crossing.}

\eqnote

As the user moves, the Positioning subsystem will continue updating their estimated location. If the system detects that the user has moved away from the planned route, it will attempt to calculate another accessible path. However, if the location information is unreliable or another suitable route cannot be found, the application will notify the user instead of giving directions that may not be correct.

{\color{red}\noindent\textbf{Note from Dr.\ Sekmen.} Graph loading and graph validation (R2) are not described yet. Also choose one search algorithm, and give the distance at which the user is considered off the route.\par}

% .....................................................................
\subsection[Accessible Interface and Guidance]{\upd{Accessible Interface and Guidance}}

\begin{figure}[htbp]
	\centering
	\includegraphics[width=0.9\textwidth]{./figures/Sub_Guidance.png}
	\caption[Accessible Interface and Guidance and its components.]{\upd{Accessible Interface and Guidance and its components.}}
	\label{fig:sub-guide}
\end{figure}

\fignote{fig:sub-guide}

The \upd{Accessible Interface and Guidance} subsystem will control how users interact with the application. Since the system is designed for blind and low-vision pedestrians, our main focus is to make navigation possible without requiring users to constantly look at their phones.

The iPhone application will allow users to select a destination, start or stop navigation, repeat instructions, and request their current location. We will also design the interface to work with Apple's VoiceOver accessibility feature, which reads information on the screen aloud. Supported voice commands will give users another way to perform common actions without using the touchscreen.

During navigation, this subsystem will receive information about the planned route, upcoming turns, and the user's estimated location. It will use this information to decide when to give directions. For example, when a user is approaching an intersection, the application will announce the upcoming turn before they reach it. This should give the user enough time to prepare for the next direction.

The system will use the iPhone's text-to-speech capabilities to provide spoken instructions. Vibration alerts will also be used for events such as turns, route deviations, warnings, and arrival at the destination. Users will be able to repeat instructions or request an update about their estimated location when needed.

We also plan to explore Apple Watch integration as an optional feature. If implemented, the watch could provide vibration alerts and simple controls, such as repeating an instruction or stopping navigation. The iPhone will still handle the main navigation processing, so users will not need an Apple Watch to use the system.

{\color{red}\noindent\textbf{Note from Dr.\ Sekmen.} Give the distance before a turn at which the instruction is spoken, and list the vibration pattern for each event in a small table.\par}

% .....................................................................
\subsection[Landmark Recognition]{\upd{Landmark Recognition}}

\begin{figure}[htbp]
	\centering
	\includegraphics[width=0.9\textwidth]{./figures/Sub_LandmarkRecognition.png}
	\caption[Landmark Recognition and its components.]{\upd{Landmark Recognition and its components.}}
	\label{fig:sub-recog}
\end{figure}

\fignote{fig:sub-recog}

For landmark recognition, camera images will be processed using the trained Core ML model stored in the application. The model will attempt to identify selected landmarks, such as building signs or entrances. When a landmark is detected, the system will check the model's confidence level and compare the result with the user's estimated location.

For example, if the camera identifies a building sign, the application will check whether that building is located near the user. If the identification is reliable and matches the location data, the application can announce the landmark to the user. This extra check is important because the model may sometimes recognize an object incorrectly.

% .....................................................................
\subsection[Glasses Connection]{\upd{Glasses Connection}}

\begin{figure}[htbp]
	\centering
	\includegraphics[width=0.9\textwidth]{./figures/Sub_GlassesConnection.png}
	\caption[Glasses Connection and its components.]{\upd{Glasses Connection and its components.}}
	\label{fig:sub-glasses}
\end{figure}

\fignote{fig:sub-glasses}

The Smart Glasses Integration and Landmark Recognition subsystems will provide additional ways for users to interact with the system and receive information about their surroundings. We plan to use compatible Meta smart glasses because they include a camera, microphone, and speakers. The iPhone will remain responsible for processing navigation information.

The Glasses Connection subsystem will manage communication between the smart glasses and the iPhone using supported Meta developer tools. Depending on the features available to developers, the glasses may allow the application to receive camera images, accept voice input, and play navigation instructions through their speakers. We will need to confirm which of these features are supported during implementation.

The Glasses Connection subsystem will also monitor whether the smart glasses are connected. If the connection is lost, the user will still be able to continue navigation using the iPhone. Where supported, the iPhone's camera, microphone, and speakers can provide some of the same functions. This will allow the main navigation features to work without depending entirely on the smart glasses.

{\color{red}\noindent\textbf{Note from Dr.\ Sekmen.} Confirm now which Meta toolkit you will use and which features it gives you (camera frames, microphone, speaker). If camera frames are not available, say that Landmark Recognition will use the phone camera.\par}

% .....................................................................
\subsection[Testing and Replay]{\upd{Testing and Replay}}

\begin{figure}[htbp]
	\centering
	\includegraphics[width=0.9\textwidth]{./figures/Sub_TestingReplay.png}
	\caption[Testing and Replay and its components.]{\upd{Testing and Replay and its components.}}
	\label{fig:sub-test}
\end{figure}

\fignote{fig:sub-test}

{\color{red}\noindent\textbf{Note from Dr.\ Sekmen.} This subsystem (R10 to R12) is in Figure~\ref{fig:architecture} but is not described in your chapter. Write this subsection. Explain how a test walk is recorded, how the recording is replayed through the app, how the replayed events are compared with the expected ones, how the supervised walking trials are run and what is measured, and how the accessibility review is done.\par}
